\section{Introduction}
Drip (trickle) irrigation delivers water at a few litres per hour to the root zone through emitters on small tubes. The design follows the water from the emitter backwards: the emitter sets the discharge, the lateral the pressure variation, the submain and main the pressure at the head, and the filter and pump the energy. A design is good when the emission uniformity is above about 90 per cent.

\section{Types of drip systems}
Drip is of three kinds: \emph{point-source} emitters (drippers on the line or in-line), \emph{line-source} (dripperline with continuous emitters) and \emph{subsurface} (SDI, buried). Emitters may be \emph{pressure compensating} (the discharge is constant over a range of pressures, $x\approx0$) or \emph{non-compensating} ($x$ from 0.5 to 1).

\section{Emitter hydraulics}
The discharge of an emitter is $q=K_dH^{x}$. A turbulent-flow emitter has $x=0.5$: with $q_0=1.6$ L\,h$^{-1}$ at $H_0=10$ m, a loss of 1.5 m of head gives $q=1.48$ L\,h$^{-1}$, a fall of 7.8 per cent.

\section{Lateral design}
The Hazen--Williams head loss in a pipe is $h_f=10.67\,L\,Q^{1.852}/(C^{1.852}D^{4.871})$ with $Q$ in m$^3$\,s$^{-1}$ and $D$ in m. For a lateral with outlets the loss is multiplied by the Christiansen factor $F=\frac{1}{m+1}+\frac{1}{2N}+\frac{\sqrt{m-1}}{6N^2}$, which for $m=1.852$ and $N=200$ outlets is 0.353. The lateral of 100 m with 200 emitters carries 0.089 L\,s$^{-1}$.

\section{Emission uniformity}
The \kw{emission uniformity} is $EU=100(1-1.27\,CV/\sqrt{n_e})\,q_{min}/q_{avg}$. It combines the manufacturing variation $CV$ (here 0.03) and the pressure variation along the lateral. An EU above 90 per cent is excellent, 80--90 acceptable and below 70 unacceptable.

